% =============================================================================
% Lane WB — Verification & honesty policy.
% Every claim grounded in: docs/architecture.md, docs/performance_model.md,
% scripts/hw_acceptance.sh, .omo/plans/tensixflow.md (guardrails),
% .omo/evidence/task-41-tensixflow.json, .omo/evidence/task-F3-tensixflow.log,
% .omo/start-work/ledger.jsonl. No external citations (contract: cite-free).
% =============================================================================
\section{Verification and Honesty Policy}\label{sec:policy}

TensixFlow's acceptance policy is built around one principle: \emph{two different
claims must never be conflated}. The first claim --- numerical correctness --- is
validated at the schedule-IR level by a functional simulator. The second claim ---
the emitted TT-Metalium kernel sources --- is validated structurally only, because
no Tenstorrent device exists anywhere in this campaign's loop. This section states
the honesty policy verbatim, describes the layered non-vacuity defense that keeps
the numeric bar from becoming meaningless, defines the disjunctive acceptance rule
for the hardware and upstream waves, and lists the scope guardrails that bound the
whole project.

\subsection{Honesty policy: generated kernels are structurally validated only}
\label{sec:policy-honesty}

The project's own documentation carries a disclaimer that this report quotes
verbatim rather than paraphrases, from \path{docs/architecture.md}
(\S\emph{Honesty disclaimer}):

\begin{quote}
``Nothing in this project executes the generated C++ kernels on hardware. There is
no Tenstorrent device in the loop, and no claim of on-device behavior anywhere.''

\smallskip
``In short: the schedule is numerically validated; the generated kernels are
structurally aligned and compile-checked. The two are different claims, and this
project keeps them separate on purpose.''
\end{quote}

Concretely, the generated kernels (\path{host.cpp}, \path{reader.cpp},
\path{compute.cpp}, \path{writer.cpp}) pass exactly two structural layers, both
executed by the test suite and never by prose:

\begin{enumerate}[leftmargin=1.6em]
  \item \textbf{Golden FileCheck patterns against the current tt-metal API}
    (pinned \texttt{@df36cd6}): the patterns assert current-API tokens --- the
    single \texttt{CreateKernel} with a \texttt{variant<ComputeConfig,\
    DataMovementConfig>}, \texttt{tt::CBIndex::c\_N} circular buffers, and the
    \texttt{distributed::Enqueue*} mesh API.
  \item \textbf{A compile-only \texttt{-fsyntax-only} check} of every generated
    kernel against signature-only tt-metal stub headers pinned to
    \texttt{@df36cd6}. The stubs carry no behavior; they exist so that an API
    drift becomes a compile error rather than a silent mismatch.
\end{enumerate}

The check is made non-vacuous by a deliberately failing negative fixture: a kernel
that calls the \emph{removed} legacy \texttt{EnqueueProgram} API must be rejected
by the stub set. The F3 hands-on QA run logged the pair explicitly: the current-API
baseline compiled, the legacy fixture failed with \texttt{EnqueueProgram}
undeclared (exit~1), and all four generated kernel files compiled
\texttt{-fsyntax-only} (\path{.omo/evidence/task-F3-tensixflow.log}, step~8).
Numeric correctness, in contrast, is asserted only at the schedule-IR level, as
detailed in Section~\ref{sec:policy-nonvacuity} and exercised end-to-end in
Section~\ref{sec:orchestration-finalwave}.

\textbf{Policy rule for numbers.} Every numeric claim in this report and in the
repository (PCC values, trip-wires, saved bytes, test counts) must originate from a
running test or harness invocation recorded in \path{.omo/evidence/}, never from
prose. This is the plan's own guardrail: ``No hand-waved correctness: every numeric
claim (PCC, bytes-saved) must come from a running test/harness, never prose.''

\subsection{Layered non-vacuity defense: keeping the PCC bar honest}
\label{sec:policy-nonvacuity}

A similarity metric is only meaningful if it can fail for the right reasons. The
similarity bar here --- Pearson correlation coefficient (PCC) between simulator
output and the NumPy golden --- could silently degenerate in two directions: a
\emph{corrupted} implementation that happens to stay above the floor, or a
\emph{vacuous} implementation where both sides compute identically so the metric
reads a meaningless perfect score. The campaign defends against both with a
layered design, and the layers are ordered from structural to empirical.

\paragraph{Rounding contract (why divergence is real).}
The two compared paths are \emph{deliberately} at different precision. The
functional simulator is the bf16 spatial model: inputs rounded to bf16, fp32
accumulation within a tile, each output tile rounded back to bf16, and fused
intermediates stored bf16 between stages. The golden reference is fp32-exact
throughout the composed graph, with no intermediate bf16 rounding. PCC between the
two therefore measures the real bf16-spatial-versus-exact cost of fusion plus tiled
execution --- and it structurally \emph{cannot} be identically~1.0.

\paragraph{Primary layer --- structural assertion + bit difference.}
Every harness run records three structural facts: the golden holds fused
intermediates in fp32, the simulator output is bf16, and the golden output dtype is
fp32 (fields \texttt{structural.*} in
\path{.omo/evidence/task-41-tensixflow.json}, true on all 16 runs). On top of that,
the harness asserts the bit-difference check
$\max|\text{sim}-\text{golden}| > 0$: if the two outputs were ever bit-identical
(the both-bf16-identical vacuous case), the run would fail deterministically. In
the recorded runs the gap is far from degenerate:
$\max|\text{sim}-\text{golden}|$ spans $2.77\times10^{-2}$ to
$4.13\times10^{-2}$ across all 16 runs of both targets
(\path{.omo/evidence/task-F3-tensixflow.log}, step~7).

\paragraph{Secondary layer --- empirically calibrated upper trip-wire.}
The upper band is \emph{not} a hard-coded constant. The trip-wire is calibrated
from the running implementation on the four documented seeds
(\texttt{0x20260727}, \texttt{0x5eed2026}, \texttt{0xc0ffee}, \texttt{0x1352777}):
take the measured good-case PCC and set the trip-wire strictly above it but
strictly below~1.0 --- in practice the midpoint between the good-case PCC
and~1.0. Hard-coding a value such as $0.9999$ was explicitly rejected: a clean
software bf16 model can exceed it, and the gate would false-fail correct code. As
recorded in the metrics JSON: for \texttt{rmsnorm\_matmul} the good-case PCC is
$0.9999940843$ and the trip-wire $0.9999970422$; for \texttt{mlp\_block},
$0.9999906926$ and $0.9999953463$. The lower bar $0.9875$ is a \textbf{fixed}
correctness floor: it is never lowered, and divergence is never injected to land
inside the band (\path{docs/performance_model.md}).

Table~\ref{tab:policy-nonvacuity} gives the full recorded picture, and
Figure~\ref{fig:policy-pcc-band} shows the pass band geometrically, including where
the two degenerate failure modes land.

\begin{table}[t]
  \centering
  \footnotesize
  \caption{Recorded PCC band per target (16 runs total: 2 targets $\times$ 2
    shapes $\times$ 4 seeds; source \texttt{.omo/evidence/task-41-tensixflow.json}).
    The trip-wire is the midpoint between the good-case PCC and 1.0, so it is
    recalibrated whenever the implementation's good case moves.}
  \label{tab:policy-nonvacuity}
  \begin{tabular}{@{}l r r r r@{}}
    \toprule
    Target & Measured PCC range & Good-case PCC & Trip-wire & Floor \\
    \midrule
    \texttt{rmsnorm\_matmul} & $0.999993159$--$0.999994084$ & $0.999994084$ & $0.999997042$ & $0.9875$ \\
    \texttt{mlp\_block}      & $0.999988845$--$0.999990693$ & $0.999990693$ & $0.999995346$ & $0.9875$ \\
    \bottomrule
  \end{tabular}
\end{table}

\begin{figure}[t]
  \centering
  \begin{tikzpicture}[
    >=Stealth,
    zonefail/.style={fill=tensixamber!16},
    zonepass/.style={fill=tensixteal!12},
    boundary/.style={line width=1pt, tensixnavy},
    trip/.style={line width=1pt, tensixamber, dashed},
    vlab/.style={font=\footnotesize, align=right},
  ]
  % --- zone fills: x band from 0.4 to 7.8 ------------------------------------
  \fill[zonefail] (0.4,0.2) rectangle (7.8,2.0);
  \fill[zonepass] (0.4,2.0) rectangle (7.8,6.8);
  \fill[zonefail] (0.4,6.8) rectangle (7.8,8.6);
  % --- axis (annotated, NOT linearly scaled) ---------------------------------
  \draw[->, line width=0.8pt, tensixnavy] (0,-0.1) -- (0,9.1)
    node[above, font=\footnotesize\bfseries, align=center] {PCC\\[-1pt]\scriptsize(annotated, not to scale)};
  % --- boundaries + exact values ---------------------------------------------
  \draw[boundary] (0.4,2.0) -- (7.8,2.0);
  \node[vlab] at (0.15,2.0) {\texttt{0.9875}\\[-1pt]\scriptsize fixed floor};
  \draw[trip] (0.4,6.8) -- (7.8,6.8);
  \node[vlab] at (0.15,6.8) {\texttt{0.999995346}\\[-1pt]\texttt{0.999997042}\\[-1pt]\scriptsize empirical trip-wire\\[-1pt]\scriptsize(mlp / rmsnorm)};
  \draw[boundary] (0.4,8.6) -- (7.8,8.6);
  \node[vlab] at (0.15,8.6) {\texttt{1.0}};
  % --- zone texts --------------------------------------------------------------
  \node[font=\footnotesize, align=center, text width=6.6cm] at (4.1,1.1)
    {\textbf{FAIL} --- corrupted or regressed run:\\ PCC below the fixed floor $0.9875$};
  \node[font=\footnotesize, align=center, text width=6.6cm] at (4.1,7.7)
    {\textbf{FAIL} --- vacuous bit-identical run:\\ $\mathrm{PCC}=1.0$ exceeds the trip-wire};
  % --- measured cluster ---------------------------------------------------------
  \node[draw=tensixteal, fill=tensixteal!25, rounded corners=2pt,
        font=\footnotesize, align=center, text width=6.2cm, inner sep=4pt]
    (cluster) at (4.1,4.3)
    {measured good-case cluster: 16 runs, 2 targets\\
     \texttt{0.99998885}\ldots\texttt{0.99999408}};
  % --- vacuous drift arrow --------------------------------------------------------
  \draw[->, dashed, tensixamber, line width=0.7pt] (cluster.north) -- (4.1,6.8)
    node[midway, right=2pt, font=\scriptsize, align=left]
      {if both paths rounded\\ identically (vacuous),\\ PCC drifts to $1.0$};
  % --- pass-band brace ---------------------------------------------------------------
  \draw[decorate, decoration={brace, amplitude=5pt, mirror}, tensixnavy, line width=0.7pt]
    (8.0,2.0) -- (8.0,6.8)
    node[midway, right=8pt, font=\footnotesize, align=left]
      {pass band:\\ floor $\le$ PCC $<$ trip-wire\\[2pt]
       \scriptsize both degenerate cases land\\ \scriptsize outside, for different reasons};
  \end{tikzpicture}
  \caption{The non-vacuity pass band on the PCC metric. The corrupted-run failure
    mode (below the fixed floor $0.9875$) and the vacuous-run failure mode
    ($\mathrm{PCC}=1.0$, bit-identical outputs) sit on opposite sides of the pass
    band and are caught by different layers: the floor and the empirically
    calibrated trip-wire respectively, with the structural fp32-golden/bf16-simulator
    assertion and the $\max|\text{sim}-\text{golden}|>0$ check underneath both.
    Axis positions are schematic; all annotated values are exact, from
    \texttt{.omo/evidence/task-41-tensixflow.json}.}
  \label{fig:policy-pcc-band}
\end{figure}

The harness encodes each layer as an executable test, all green in the final
hands-on QA run (11/11 in \path{test_correctness.py}):
\path{test_pcc_above_floor},
\path{test_max_abs_diff_positive_not_bit_identical},
\path{test_pcc_below_empirical_tripwire},
\path{test_tripwire_calibrated_above_good_case}, and
\path{test_structural_fp32_golden_bf16_simulator}, each parametrized over
both targets.

\subsection{Disjunctive acceptance policy}
\label{sec:policy-disjunctive}

Two acceptance gates depend on resources the campaign does not control: a
Tenstorrent device (not owned) and an upstream maintainer's merge decision (not
owned). Both are therefore phrased \emph{disjunctively} --- each passes on a
documented fallback as well as on the full path --- with one strict structural
requirement on how the disjunction is written.

\paragraph{Hardware wave: explicit \texttt{if/else}, never \texttt{\&\&\ldots||}.}
The acceptance harness \path{scripts/hw_acceptance.sh} probes the device node and
runs exactly one branch:

\begin{lstlisting}[basicstyle=\ttfamily\scriptsize, frame=single, caption={Disjunctive hardware acceptance shape --- verbatim skeleton from \texttt{scripts/hw\_acceptance.sh} (\texttt{TT\_DEVICE\_PATH} defaults to \texttt{/dev/tenstorrent}).}, label={lst:policy-hwaccept}]
if [ -e "${TT_DEVICE_PATH}" ]; then
  # ON-DEVICE branch (todo 55): compile + run generated kernels,
  # validate PCC >= 0.9875 vs the NumPy golden.
  bash scripts/hw_run_on_device.sh
else
  # NO-DEVICE branch (todo 56): runbook dry-run + compile-only check.
  bash scripts/hw_bringup_dryrun.sh \
    && cmake --build build --target check-generated-kernels-compile
fi
\end{lstlisting}

The \texttt{if/else} is deliberate. An \texttt{on-device \&\&\ldots\ ||\
no-device} form would let a \emph{present-but-failing} device run fall through to
the passing no-device branch and be silently masked; with \texttt{if/else} exactly
one branch runs and the script's exit status is that branch's status, unmasked. The
on-device arm fails honestly even in absence of its entry point: if
\texttt{hw\_run\_on\_device.sh} is missing or non-executable the script exits~1
and refuses to fall through. On this campaign's CPU-only box the no-device arm is
the active one and did real work: bring-up dry-run exit~0 plus
\texttt{check-generated-kernels-compile} exit~0 (baseline compiles, legacy fixture
rejected, four generated kernels compile) --- see the final-wave QA in
Section~\ref{sec:orchestration-finalwave}. The on-device branch itself (todo~55)
was skipped and is recorded as such in the ledger, because no device is owned.

\paragraph{Upstream wave: submitted, never ``merged''.}
The upstream contribution gate accepts either an open pull request with the
\emph{contributor's own} checks green, or a prepared \path{0001-*.patch} plus an
issue URL. It never depends on the maintainers' CI and never claims a merge, since
neither is owned by this campaign. In the event, the campaign owner issued an
explicit do-not-subscribe/do-not-submit override, so todo~T60 landed as
\emph{documented-only}: the scouted target, the contribution, and the submission
procedure are committed under \path{upstream/}, and the ledger records
``T60 skipped --- submission process documented only.''

\subsection{Scope guardrails as policy}
\label{sec:policy-guardrails}

The plan's ``Must NOT have'' list is enforced as acceptance policy, not as advice:
each stretch surface has a binding guardrail and an executable or auditable check.
Table~\ref{tab:policy-guardrails} reproduces them together with how each was
verified during the campaign (final-wave evidence cited where applicable; see
Section~\ref{sec:orchestration-finalwave}).

\begin{table}[t]
  \centering
  \footnotesize
  \caption{Scope guardrails binding the campaign, and the check that enforced each.}
  \label{tab:policy-guardrails}
  \begin{tabular}{@{}p{2.7cm} p{6.1cm} p{5.2cm}@{}}
    \toprule
    Surface & Guardrail (binding) & Enforced / verified by \\
    \midrule
    Python DSL & Thin MLIR-text builder only; \emph{no} tracing JIT, no AST
      introspection & Canonicalized-IR equivalence test vs the handwritten inputs
      (ignoring SSA names); F4 import audit: no \texttt{ast}/\texttt{inspect} \\
    TT-MLIR compat & TTIR \emph{emitter} only; the other four dialects are prose
      mappings & \texttt{--tensixflow-emit-ttir} lit test; mapping doc covers
      TTIR/TTNN/D2M/TTKernel/TTMetal in prose \\
    Remote-box script & \texttt{bash -n} syntax check + docs only; never live
      provisioning, no cloud spend & \texttt{bash -n scripts/remote-box.sh}
      exit~0; script never executed; pricing prose dated ``as-of 2026-07-21'' \\
    Benchmarks & Timing loops, \emph{not} a framework & No benchmark-framework
      dependency (\texttt{find\_package(benchmark)} absent); each binary prints an
      explicit static-estimate disclaimer \\
    HTML visualizer & Static, offline, zero external references & Playwright load
      with all network requests aborted: 0 requests, 0 console errors, 6/6 section
      anchors; 0 \texttt{http(s)} occurrences in the source \\
    Owner constraints & nix/devenv shell binding the whole build; no NVIDIA GPU;
      no PR submission (T60 documented-only) & Every gate runs under
      \texttt{nix develop}; F3 probes: no \path{/dev/tenstorrent}, no NVIDIA in
      \texttt{lspci}; F4 audit: nothing submitted/forked/merged \\
    \bottomrule
  \end{tabular}
\end{table}

Two further guardrails shape everything downstream: exactly \emph{one} modeled
device (Wormhole n150-class: 8$\times$8 grid, 32$\times$32 tiles, bf16/fp32,
1~MiB modeled L1 per core) --- no multi-target support --- and no NVIDIA/CUDA path
at all, since this is a spatial-accelerator compiler and the campaign owns no
NVIDIA GPU. Both were re-audited adversarially by the F4 scope-fidelity review
(repo-wide greps for \texttt{cuda|nvidia}, \texttt{grayskull|blackhole|n300}, and
multi-device tokens came back clean except a single sanctioned
absence-statement in a test docstring). Hardware specifics are covered in
Section~\ref{sec:hardware}.
